% section: 式の変形

\section{式の変形}\label{節：式の変形}
  \noindent\textbf{【本編とのつながり】}\par
  同じ値を保ちながら既知の形を作る操作を確認する.定数の調整は\sectionref*{節：特性方程式型},\bookterm{partial}{部分分数分解}は\sectionref*{節：階差型}で和を計算する際に使う.

  \subsection{同じ数を足して引く}\label{節：+a-aをする}
    $x^2+6x+5$を,一つの2乗を使って表したい.
    $(x+3)^2=x^2+6x+9$なので,$x^2+6x$の部分に$9$を足せば平方になる.値を変えないため,同じ$9$を引く.
    \[
      \begin{aligned}
        x^2+6x+5&=x^2+6x+9-9+5\\
        &=(x+3)^2-4.
      \end{aligned}
    \]
    値を変えずに,$x$を含む部分を一つの2乗にまとめられた.
    この操作を\bookterm{square}{平方完成}という.
    一般の二次式でも,$a\ne0$として同じように進める.
    \begin{align*}
      &ax^2+bx+c\\
      &\qquad=a\left(x^2+\frac{b}{a}x\right) +c\\
      &\qquad=a\left(x^2+\frac{b}{a}x+\frac{b^2}{4a^2}-\frac{b^2}{4a^2}\right)+c \booktag{square-adjust}{☆}\\
      &\qquad=a\left(x+\frac{b}{2a}\right)^2-\frac{b^2}{4a}+c
    \end{align*}
    と,$x$に関する項を1つにまとめられる.
    ここで注目してほしいのは\bookeqbackref{square-adjust}{☆}部分の式変形である.
    これが意味するのは,無理やりに2乗の形,すなわち平方を作り上げよう,ということだ.
    これにより$x$に関する項は定数とともに括弧の中に閉じ込められ,よりシンプルに式が記述できる.
    二次関数の最大・最小問題や軸の決定などにおいても\bookterm{square}{平方完成}は有用である.

    ここから学ぶべきことは,定数項のやり取りによって式を無理やりわかりやすい形,知っている形,扱いやすい形に変形できるということだ.

  \subsection{部分分数分解}\label{節：部分分数分解}
    今度は,一つの分数を簡単な分数の和に分けてみよう.
    例えば,$\dfrac1{x+1}-\dfrac1{x+2}$を通分すると,
    \[
      \begin{aligned}
        \frac1{x+1}-\frac1{x+2}
          &=\frac{(x+2)-(x+1)}{(x+1)(x+2)}\\
          &=\frac1{(x+1)(x+2)}
      \end{aligned}
    \]
    となる.ここでは$x\ne-1,-2$とする.
    この等式を逆向きに使えば,分母が二次式の分数を,分母が一次式の二つの分数へ分けられる.
    $x$を一つずつ増やして和を取ると,隣の項と打ち消し合う形でもある.
    このような操作を\bookterm{partial}{部分分数分解}という.

    一般に
    \[
      \frac{1}{(x+a)(x+b)}
    \]
    でも,分母を一次式に分けることを考える.
    以下では$x\ne-a,-b$とし,$a\ne b$の場合を扱う.
    \[
      \frac{1}{(x+a)(x+b)}=\frac{\alpha}{x+a}+\frac{\beta}{x+b}
    \]
    分母が零にならない\textbf{すべての$x$で一致する}ように,$\alpha,\beta$を決める.
    右辺を通分すると,
    \begin{align*}
    \frac{1}{(x+a)(x+b)}
      &= \frac{\alpha}{x+a}+\frac{\beta}{x+b}\\
      &= \frac{(\alpha+\beta)x+\alpha b+\beta a}
      {(x+a)(x+b)}
    \end{align*}
    分子が常に$1$になるには,$x$の係数が$0$,定数項が$1$であればよい.
    \[
    \therefore\quad
    \begin{cases}
    \alpha+\beta=0,\\
    \alpha b+\beta a=1
    \end{cases}
    \]
    これを解くと,\,$\beta=-\alpha$であるから,
    \[
      \alpha b-\alpha a=1
      \iff \alpha(b-a)=1
    \]
    したがって,\,$a\ne b$であれば,
    \[
      \alpha=\frac{1}{b-a},
      \qquad
      \beta=-\frac{1}{b-a}
    \]
    が得られる.従って
    \[
      \frac1{(x+a)(x+b)}
      =\frac1{b-a}\left(\frac1{x+a}-\frac1{x+b}\right).
    \]
    $a=b$のときは分母に同じ因子が二つ現れるため,この形の分解はできない.
    3次の場合も,\,$p,q,r\ne0$で,分母の一次因子$px+a,qx+b,rx+c$の根
    \[
      -\frac{a}{p},\qquad -\frac{b}{q},\qquad -\frac{c}{r}
    \]
    が互いに異なるときは,
\BookMathLayout{\[
      \frac{1}{(px+a)(qx+b)(rx+c)}=\frac{\alpha}{px+a}+\frac{\beta}{qx+b}+\frac{\gamma}{rx+c}
    \]}{
\[
\begin{aligned}
\frac{1}{(px+a)(qx+b)(rx+c)}
&=\frac{\alpha}{px+a}+\frac{\beta}{qx+b}\\
&\quad+\frac{\gamma}{rx+c}
\end{aligned}
\]
}
    と置き,係数$\alpha,\beta,\gamma$を定めることで分解できる.
    言い換えれば,これらの一次因子が互いに定数倍でないことが条件である.
    重複因子があるときは,その重複度に応じて例えば$\dfrac{A}{x+a}+\dfrac{B}{(x+a)^2}$のような項も必要になる.
